
\section{Введение}
Метод Хигучи используется для характеристики масштабных свойств временных рядов~\cite{higuchi1988}. Конечновыборочная оценка HFD зависит от длины наблюдения и диапазона масштабов; математическая интерпретация оценок требует осторожности~\cite{liehr2020,wanliss2022}. Для зависимых наблюдений разработаны методы блочной повторной выборки~\cite{kunsch1989,politis1994,hall1995}, но их общие асимптотические результаты не гарантируют номинального покрытия конкретного нелинейного функционала на коротких отрезках.
Цель работы --- эмпирически изучить распределение статистики Хигучи $T_{N,16}$ и покрытие доверительных интервалов для ее математического ожидания при фиксированных стационарных моделях. Ни новый bootstrap-алгоритм, ни теоретические гарантии для реальных EEG здесь не предлагаются.

\section{Материалы и методы}
\subsection{Определение статистики}
Для ряда $X_1,\ldots,X_N$ и $1\leq k\leq 16$ положим
\begin{equation}
 n_{a,k}=\lfloor(N-a)/k\rfloor,\quad 1\leq a\leq k,
\end{equation}
\begin{equation}
 L_a(k)=\frac{N-1}{n_{a,k}k^2}
 \sum_{j=1}^{n_{a,k}}\left|X_{a+jk}-X_{a+(j-1)k}\right|,
\end{equation}
\begin{equation}
 L(k)=\frac1k\sum_{a=1}^{k}L_a(k).
\end{equation}
Пусть $u_k=\log(1/k)$, $v_k=\log L(k)$ и $\bar u=16^{-1}\sum_{k=1}^{16}u_k$. Статистика Хигучи --- наклон регрессии наименьших квадратов:
\begin{equation}\label{eq:hfd}
 T_{N,16}=\frac{\sum_{k=1}^{16}(u_k-\bar u)v_k}
 {\sum_{k=1}^{16}(u_k-\bar u)^2}.
\end{equation}
Если хотя бы один $L(k)$ не является конечным положительным числом, $T_{N,16}$ не определена. Целевой параметр при модели $P$:
\begin{equation}\label{eq:theta}
 \theta_{N,16}(P)=\mathbb{E}_P[T_{N,16}].
\end{equation}
Это конечновыборочное математическое ожидание статистики, а не теоретическая фрактальная размерность процесса или прогноз HFD нового фрагмента.

\subsection{Стационарные модели и дизайн}
Рассматриваются нормированные к единичной дисперсии гауссовские модели:
\begin{equation}
 \mathrm{M1}: X_t=0.3X_{t-1}+\varepsilon_t,\quad
 \varepsilon_t\sim N(0,1-0.3^2);
\end{equation}
\begin{equation}
 \mathrm{M2}: X_t=0.8X_{t-1}+\varepsilon_t,\quad
 \varepsilon_t\sim N(0,1-0.8^2);
\end{equation}
\begin{equation}
 \mathrm{M3}: X_t=\frac{\varepsilon_t+0.6\varepsilon_{t-1}}
 {\sqrt{1+0.6^2}},\quad \varepsilon_t\sim N(0,1).
\end{equation}
В M1 и M2 $X_0\sim N(0,1)$ задается стационарно; инновации независимы. Использованы длины $N\in\{384,768\}$, $K=16$, $R=200$ исходных реализаций для каждой модели и длины, $B=199$ повторных выборок на метод и исходную реализацию, по $20\,000$ независимых эталонных рядов для оценки $\theta_{N,16}$; базовое значение генератора seed равно 20261029. Подробная реализация и конфигурация доступны в экспериментальном архиве.

\subsection{Bootstrap и интервалы}
Применены moving block bootstrap с перекрывающимися блоками длины 32 и 96 (MBB32, MBB96); stationary bootstrap со средней длиной блока 32 и 96 (SB32, SB96); независимый IID bootstrap используется как контроль, игнорирующий временную зависимость. В stationary bootstrap начало нового блока выбирается с вероятностью $1/b$, а внутри блока индексы продолжаются циклически.
Пусть $q^*_\alpha$ --- квантиль распределения $T^*$ при фиксированном исходном ряде. Основной предзаданный 95\%-й интервал имеет вид
\begin{equation}\label{eq:basic}
 I_{\mathrm{basic}}=[2T-q^*_{0.975},\,2T-q^*_{0.025}].
\end{equation}
Квантили получены методом линейной интерполяции \texttt{numpy.quantile} (\texttt{method=linear}); число bootstrap-повторений $B=199$. Только дополнительно и диагностически по сохраненным тем же квантилям рассмотрен
\begin{equation}\label{eq:perc}
 I_{\mathrm{perc}}=[q^*_{0.025},\,q^*_{0.975}].
\end{equation}
Оба варианта проверены относительно параметра~\eqref{eq:theta}, а не теоретической размерности или будущей реализации.

Эмпирическое покрытие вычислялось как $\hat p=R^{-1}\sum_{r=1}^R\mathbf 1\{\hat\theta_{\mathrm{ref}}\in I_r\}$. Monte Carlo standard error равна $\sqrt{\hat p(1-\hat p)/R}$; для частоты покрытия также вычислены интервалы Уилсона. Последние обусловлены фиксированным численным эталоном $\hat\theta_{\mathrm{ref}}$ и не учитывают его собственной ошибки. Дополнительные критерии исходного протокола включают смещение и RMSE оценки дисперсии; один лишь порог покрытия не является полным критерием успеха. Сравнения интервалов выполнялись попарно для одинаковых исходных реализаций.

\section{Результаты}
\subsection{Покрытие исходных basic-интервалов}
Таблица~\ref{tab:coverage} воспроизводит все 30 проверенных значений исходных CSV. Каждая доля основана на $R=200$ Monte Carlo-реализациях.
\begin{table}[!t]
\caption{Эмпирическое покрытие basic CI (доля)}\label{tab:coverage}
\centering\scriptsize
\begin{tabular}{llrrrrr}
\toprule
Модель&$N$&MBB32&MBB96&SB32&SB96&IID\\
\midrule
M1&384&.855&.815&.855&.810&.175\\
M2&384&.800&.780&.780&.785&.000\\
M3&384&.895&.780&.855&.805&.075\\
M1&768&.910&.910&.890&.880&.070\\
M2&768&.800&.855&.780&.845&.000\\
M3&768&.860&.875&.885&.845&.005\\
\bottomrule
\end{tabular}
\end{table}
При $N=384$ ни одна из четырех блочных конфигураций не обеспечила заранее заданное покрытие не ниже 90\% одновременно на M1--M3. Для M2, $N=384$, MBB32 покрытие составило $0.800$, MCSE $=0.0283$ (95\%-й интервал Уилсона приблизительно $[0.739;0.849]$). Доля 0.93 для диагностического percentile-интервала имеет MCSE $\approx0.0180$. Полная таблица частот, интервалов Уилсона и ширин хранится в воспроизводимых CSV.

\subsection{Парные сравнения и чувствительность}
Для M2, $N=384$, парная разность покрытия MBB32 минус MBB96 равна $+0.020$ (приближенный 95\%-й интервал $[-0.0238;0.0638]$), а MBB32 минус SB32 равна $+0.020$ ($[-0.0167;0.0567]$). Эти результаты не дают убедительных оснований считать MBB32 предпочтительнее указанных методов.
Для той же модели и длины MBB32 basic покрывает 80\%, а диагностический percentile --- 93\%: разность percentile минус basic составляет $+13$ процентных пунктов. Только basic покрывает эталон в 5 реализациях, только percentile --- в 31 реализации; всего 36 несовпадений. Результат получен на уже использованных квантилях и \textit{не является независимой валидацией} percentile bootstrap. При редких несовпадениях приближенные Wald-интервалы для парных разностей не имеют строгой гарантии точности.

Проверка чувствительности выполнялась без повторной генерации данных путем подстановки $\hat\theta_{\mathrm{ref}}\pm1.96\,SE(\hat\theta_{\mathrm{ref}})$ во все 30 комбинаций. Максимальный размах трех значений покрытия равен двум процентным пунктам: для M2, $N=384$, MBB32 покрытие составило соответственно 81.5\%, 80.0\%, 79.5\%. Для остальных конфигураций изменение не превосходило одного процентного пункта. Это диагностическая проверка устойчивости основного вывода, не полный учет двух уровней Monte Carlo uncertainty.

\subsection{Диагностика искусственных стыков}
В диагностике первых десяти исходных реализаций исследованы распределения приращений с лагом 1 на искусственных стыках bootstrap-блоков. Их различия зарегистрированы, однако статистика~\eqref{eq:hfd} использует приращения с лагами $1,\ldots,16$. Причинный вклад стыков в смещение статистики и недопокрытие интервалов данным экспериментом не идентифицирован.

\section{Обсуждение}
Классические результаты о блочном bootstrap обосновывают рассмотрение ресемплинга зависимых последовательностей~\cite{kunsch1989,politis1994,hall1995}, но не являются готовым утверждением о надежности коротких интервалов для нелинейной статистики Хигучи. Наблюдаемое недопокрытие не позволяет установить, какая часть вызвана смещением условного bootstrap-распределения, нарушением временной зависимости на стыках или ошибкой оценки дисперсии. Диагностическое различие basic и percentile показывает чувствительность к формуле интервала, но не доказывает общего преимущества percentile CI.

Применимость результата ограничена тремя конкретными гауссовскими стационарными моделями, двумя длинами и пятью вариантами повторной выборки. Эталонный параметр определен с помощью 20\,000 симуляций, а Wilson CI не включает его погрешность, хотя отдельная проверка чувствительности не меняет основного вывода. Полный Monte Carlo независимо не перезапускался; аудит проверил исходные числа, код, тесты и повторный анализ, но не является полной независимой репликацией. Не установлены ни общая несостоятельность блочного bootstrap, ни причинный механизм недопокрытия, ни валидность исследованных интервалов для реальных Motor Imagery EEG. Сопоставление с литературой носит ограниченный характер и не обосновывает заявление о первенстве исследования.

\section{Заключение}
Для заранее фиксированной статистики $T_{N,16}$ изучено покрытие существующих bootstrap-интервалов для конечновыборочного среднего $\theta_{N,16}$. При $N=384$ ни одна блочная конфигурация не достигла предустановленного порога покрытия во всех трех моделях. Покрытие чувствительно к построению интервала в отдельных конфигурациях, однако диагностический percentile-анализ не заменяет независимую проверку. Результат представляет ограниченное воспроизводимое методологическое сравнение, не новый bootstrap-алгоритм и не гарантию для EEG.

\begin{thebibliography}{9}
\bibitem{higuchi1988} Higuchi T. Approach to an irregular time series on the basis of the fractal theory. \textit{Physica D}. 1988. DOI: 10.1016/0167-2789(88)90081-4.
\bibitem{liehr2020} Liehr L., Massopust P. On the mathematical validity of the Higuchi method. \textit{Physica D}. 2020. arXiv:1906.10558.
\bibitem{wanliss2022} Wanliss J., Wanliss G. Efficient calculation of fractal properties via the Higuchi method. \textit{Nonlinear Dynamics}. 2022. DOI: 10.1007/s11071-022-07353-2.
\bibitem{kunsch1989} Künsch H. R. The jackknife and the bootstrap for general stationary observations. \textit{Annals of Statistics}. 1989. DOI: 10.1214/aos/1176347265.
\bibitem{politis1994} Politis D. N., Romano J. P. The stationary bootstrap. \textit{Journal of the American Statistical Association}. 1994. DOI: 10.1080/01621459.1994.10476870.
\bibitem{hall1995} Hall P., Horowitz J. L., Jing B.-Y. On blocking rules for the bootstrap with dependent data. \textit{Biometrika}. 1995. DOI: 10.1093/biomet/82.3.561.
\end{thebibliography}
